\section{Conclusions}\label{sec:conclusions}

Selecting the number of clusters on spatial data is compromised by an
assumption that the standard criteria make and that spatial data violate: that
observations are independent. We have shown that the violation is not benign.
Under spatial autocorrelation the gap statistic and the silhouette,
Calinski--Harabasz and Davies--Bouldin indices over-detect without exception,
inventing classes on featureless fields and fragmenting single land-cover
regions on real imagery.

We have proposed a remedy framed in geostatistical terms. Recast as a
calibrated test of cluster homogeneity, asking whether a region carries
multimodality in excess of spatial autocorrelation, ESS-Dip evaluates the
Hartigan dip along the locally optimal split direction and calibrates it
against the \emph{effective} sample size, estimated from a within-class
autocorrelation range obtained by local range estimation, after removing a
low-order trend whenever one is resolvable. The procedure carries an asymptotic
false-split-control property and runs in near-linear time. On a factorial
simulation with known ground truth it reaches unit specificity, and its recall
variant ESS-Dip-R attains the best balanced score of fifteen criteria at the
same specificity; on the Indian Pines, Salinas and Sentinel-2 scenes both
identify single land-cover regions as homogeneous where the standard criteria
do not. Their reliable negatives suit the family to use as a homogeneity guard
or a stopping rule for region-growing as much as to cluster-number selection,
spanning a precision--recall axis from the conservative distribution-free null
to the higher-recall Gaussian one.

The broader message is that number-of-clusters selection on autocorrelated data
must be calibrated to the information the data actually contain, not to their
nominal size, and that the signal distinguishing classes from autocorrelation
is modality, not dispersion. The effective sample size and declustering,
long-standing tools of geostatistics, supply exactly what is needed to make
a classical clustering criterion honest about spatial dependence. All code and
data required to reproduce the simulation study and the applications are
released.
